% Generated draft: requires mathematical and source review.
\documentclass[11pt]{article}
\usepackage[margin=25mm]{geometry}
\usepackage{fontspec}
\setmainfont{Latin Modern Roman}
\usepackage{amsmath,amssymb,mathtools,unicode-math}
\setmathfont{Latin Modern Math}
\usepackage{enumitem,needspace,xurl,hyperref}
\setlength{\parindent}{0pt}
\setlength{\parskip}{5pt}
\setlength{\emergencystretch}{3em}
\newcommand{\N}{\mathbb{N}}
\newcommand{\Z}{\mathbb{Z}}
\newcommand{\Q}{\mathbb{Q}}
\newcommand{\R}{\mathbb{R}}
\newcommand{\C}{\mathbb{C}}
\newcommand{\F}{\mathbb{F}}
\newcommand{\A}{\mathbb{A}}
\newcommand{\PP}{\mathbb{P}}
\newcommand{\E}{\mathbb{E}}
\newcommand{\eps}{\varepsilon}
\newcommand{\dd}{\,\mathrm{d}}
\begin{document}
\section*{UnsolvedMath ID: 30001040}
\textbf{Source problem code:} catalogue entry
\subsection{Definitions and mathematical statement}
Write \(T^2=\R^2/\Z^2\). A torus homeomorphism \(h:T^2\to T^2\) is minimal if every orbit is dense, and uniquely ergodic if it has exactly one invariant Borel probability measure \(\lambda\). The source additionally requires \(\lambda\) to be ergodic and to have full support.

The adding-machine system is read as the dyadic odometer suggested by the supplied assessment: addition of \(1\) on \(\varprojlim_r\Z/2^r\Z\), with its Haar probability measure. A two-sided Bernoulli system consists of a product probability space \((A^\Z,\nu^\Z)\) and the shift; the alphabet and probability distribution are not given by the record.

For each of these two measure-preserving systems separately, does there exist a minimal uniquely ergodic torus system \((T^2,\lambda,h)\), with the stated full-support ergodic measure, measurably conjugate to it? Measurable conjugacy means an invertible measure-preserving map between conull invariant subsets which intertwines the transformations almost everywhere. The choices of \(h\) and \(\lambda\) may be different for the two systems; the question does not require one torus transformation to realize both. The notation \(\lambda\) alone does not impose Lebesgue measure. The homeomorphism and two-sided-shift conventions explain the usual invertible reading and need confirmation if a broader meaning of topological system was intended.
\subsection{Short English statement}
Seek separate strictly ergodic models on the two-torus for the adding machine and the intended Bernoulli shift, with measurable conjugacy to their invariant probability systems. The Bernoulli parameters and any requirement that the torus measure be Lebesgue are unspecified.
\subsection{Sources}
\begin{itemize}
\item Local UnsolvedMath registry, permanent ID \(30001040\), source code \texttt{catalogue entry}, ``Torus Realizations of Adding Machines and Bernoulli Shifts.'' Registry curation: Ulam AI / UnsolvedMath Contributors, CC BY 4.0; underlying sources retain their supplied provenance.
\item \emph{Geometric Group Theory, Hyperbolic Dynamics and Symplectic Geometry} (2009), Oberwolfach Reports, pp.\ 2337--2406, DOI \url{https://doi.org/10.4171/owr/2008/42}.
\item Supplied ``Strictly ergodic models'': \url{https://arxiv.org/abs/2104.09604}.
\item Supplied literature triage, \texttt{literature\_checked\_at}: 2026-08-21: labeled open; says the standard irrational-spectrum realization theorem does not apply to the adding machine because its spectrum is dyadic, and no resolving torus realization was found. This is supplied dated curation, not independently verified; it does not separately assess the Bernoulli clause.
\end{itemize}
\end{document}
